%%%%%%%%%%%%%%%%%%%%%%%%%%%%%%%%%%%%%%%%%%%%%%%%%%%%%%%%%%%%%%%%%%%%%%%%%%%
%
% Generic template for TFC/TFM/TFG/Tesis
%
% By Javier Macías-Guarasa and contributors.
%
%%%%%%%%%%%%%%%%%%%%%%%%%%%%%%%%%%%%%%%%%%%%%%%%%%%%%%%%%%%%%%%%%%%%%%%%%%%

\chapter{Conclusiones y líneas futuras}
\label{cha:conclusiones}

\section{Conclusiones}

Con esta plantilla hemos intentado reunir en un mismo proyecto la estructura del documento académico, las cubiertas institucionales, la configuración compartida y algunos documentos complementarios que suelen acompañar a un TFG, un TFM o una tesis doctoral. También hemos querido conservar ejemplos que te permitan descubrir posibilidades de \LaTeX{} sin convertir esta guía en un manual exhaustivo.

No esperamos que utilices todo lo que contiene el repositorio. Si estás comenzando, basta con que configures tus datos, prepares los capítulos de tu trabajo y generes el PDF desde tu editor habitual. El resto de las prestaciones queda a tu disposición para cuando necesites documentos administrativos, automatización, control de cambios o ejemplos de composición más elaborados.

La plantilla no puede sustituir a la normativa de cada titulación ni a una revisión cuidadosa del resultado final. Esperamos, sin embargo, que te ahorre trabajo repetitivo y te permita dedicar más tiempo al contenido de tu proyecto. Si durante ese proceso encuentras un error o mejoras alguna de sus partes, estaremos encantados de conocer tu experiencia y de que la compartas con quienes utilicen la plantilla después.

\section{Líneas futuras}

Una plantilla de este tipo nunca puede considerarse completamente terminada: cambian las normativas, aparecen nuevos formatos institucionales y evolucionan tanto las distribuciones de \LaTeX{} como las herramientas con las que trabajamos. Entre las posibles líneas de evolución del proyecto consideramos especialmente útiles las siguientes:

\begin{itemize}
  \item Mantener actualizadas las cubiertas y los documentos administrativos conforme cambien los requisitos de las titulaciones compatibles.
  \item Ampliar el soporte a otras titulaciones e instituciones cuando dispongamos de sus modelos oficiales y de personas que puedan ayudarnos a validarlos.
  \item Simplificar progresivamente la configuración y detectar de forma más clara las combinaciones de variables incompletas o incoherentes.
  \item Mejorar la portabilidad del proceso de compilación entre GNU/Linux, Windows y macOS sin obligarte a utilizar una herramienta de construcción concreta.
  \item Incorporar comprobaciones automáticas que permitan verificar las distintas configuraciones y reducir el riesgo de que un cambio afecte a documentos que no estamos utilizando en ese momento.
  \item Seguir revisando los ejemplos para que resulten útiles, actuales y fáciles de trasladar a un documento real.
\end{itemize}

Estas líneas futuras no constituyen un plan cerrado. El rumbo de la plantilla dependerá también de los problemas reales que encuentres y de las aportaciones que recibamos. Como ha ocurrido desde sus primeras versiones, muchas de las mejoras más valiosas pueden comenzar con una pregunta, una corrección o un ejemplo compartido por uno de sus usuarios.

%%% Local Variables:
%%% TeX-master: "../book"
%%% End:
